\documentclass[a4paper,12pt]{article}
\usepackage[utf8]{inputenc}
\usepackage[T1]{fontenc}
\usepackage[french]{babel}
\usepackage{amsmath, amssymb, amsfonts}
\usepackage{array, longtable, booktabs}
\usepackage[table]{xcolor} % option table pour \rowcolors
\usepackage{geometry}
\geometry{left=1.8cm, right=1.8cm, top=2cm, bottom=2cm}
\usepackage{enumitem}
\usepackage{hyperref}
\hypersetup{colorlinks=true, linkcolor=blue, urlcolor=blue}

% Commandes pour les coches
\newcommand{\fait}{\ensuremath{\checkmark}}
\newcommand{\revoit}{\ensuremath{\circlearrowright}}
\newcommand{\casevide}{\ensuremath{\square}}

% Types de colonnes
\newcolumntype{C}[1]{>{\centering\arraybackslash}p{#1}}
\newcolumntype{L}[1]{>{\raggedright\arraybackslash}p{#1}}
\newcolumntype{R}[1]{>{\raggedleft\arraybackslash}p{#1}}

% Pour réduire les débordements dans les longtable
\setlength{\LTleft}{-0.5cm}
\setlength{\LTright}{-0.5cm}

\title{Programmation annuelle de Mathématiques -- Première technologique (tronc commun) \\
	Année 2026--2027 -- Zone C (3 h/semaine)}
\author{}
\date{Version du \today}

\begin{document}
	
	\maketitle
	
	\section*{Présentation générale}
	Ce document propose une organisation spiralaire de l’enseignement des mathématiques en première technologique (tronc commun), conforme au programme officiel. L’année comporte 31 semaines de cours effectifs (hors vacances et jours fériés). Chaque notion est introduite une première fois, puis réinvestie et approfondie dans un second temps, à l’occasion de thèmes d’étude différents. Cette approche favorise une acquisition progressive et durable des connaissances.
	
	\medskip
	\textbf{Place de Python :} La programmation est intégrée de façon transverse et évaluée dans \emph{tous} les sujets (formatives, contrôles, DM, PBL). Les compétences algorithmiques (boucles, tests, fonctions, simulation, manipulation de listes) sont développées en parallèle des mathématiques et réactivées régulièrement selon le même principe spiralé.
	
	\medskip
	\textbf{Utilisation de l’IA dans les DM :} Dans chaque devoir maison, l’élève est invité à utiliser un outil d’intelligence artificielle générative (ChatGPT, Mistral, Copilot, etc.) pour :
	\begin{itemize}
		\item obtenir des explications ou des exemples supplémentaires sur une notion,
		\item suggérer des pistes de résolution,
		\item débuguer ou améliorer un code Python,
		\item confronter différentes méthodes de résolution.
	\end{itemize}
	\emph{L’IA ne doit cependant pas rédiger la solution finale à la place de l’élève.} Le raisonnement personnel et la rédaction restent exigibles. Un \textbf{journal de bord} (questions posées, réponses obtenues, critiques et adaptations apportées) est à annexer à chaque DM. Cette pratique développe l’esprit critique et la capacité à évaluer la pertinence des réponses générées.
	
	\section{Tableau récapitulatif des contenus à enseigner}
	Le tableau ci-dessous reprend l’intégralité des attendus du programme. La colonne \textbf{« Traité »} indique si le contenu a été abordé en classe (coche \fait), et la colonne \textbf{« Reprise »} signale les notions qui feront l’objet d’un second passage (\revoit) pour consolider la compréhension.
	
	\rowcolors{2}{gray!8}{white}
	\begin{longtable}{|L{13.5cm}|C{1.5cm}|C{1.5cm}|}
		\hline
		\textbf{Thème / Contenu} & \textbf{Traité} & \textbf{Reprise} \\
		\hline
		\endhead
		\hline
		\multicolumn{3}{r}{Suite de la page suivante} \\
		\endfoot
		\hline
		\endlastfoot
		
		% --- Suites numériques ---
		\multicolumn{3}{|L{13.5cm}|}{\textbf{Suites numériques}} \\
		Différents modes de génération d’une suite (relation fonctionnelle, récurrence) & \casevide & \casevide \\
		Sens de variation d’une suite, représentation graphique (nuage de points) & \casevide & \casevide \\
		Suites arithmétiques : relation de récurrence, sens de variation, représentation & \casevide & \casevide \\
		Suites géométriques (termes strictement positifs) : relation de récurrence, sens de variation, représentation & \casevide & \casevide \\
		Reconnaître si une situation relève d’une croissance linéaire ou exponentielle & \casevide & \casevide \\
		\hline
		
		% --- Fonctions de la variable réelle ---
		\multicolumn{3}{|L{13.5cm}|}{\textbf{Fonctions – généralités et polynômes degré 2}} \\
		Taux de variation, pente de la sécante & \casevide & \casevide \\
		Fonctions monotones, lien avec le signe du taux de variation & \casevide & \casevide \\
		Fonctions polynômes de degré 2 : représentations $x \mapsto x^2$, $x^2+b$, $a(x-x_1)(x-x_2)$ & \casevide & \casevide \\
		Racines, signe, extremum, axe de symétrie d’un polynôme du second degré donné sous forme factorisée & \casevide & \casevide \\
		Résolution d’équations $x^2 = c$ & \casevide & \casevide \\
		\hline
		\multicolumn{3}{|L{13.5cm}|}{\textbf{Fonctions polynômes de degré 3}} \\
		Représentations de $x \mapsto ax^3$, $ax^3+b$ & \casevide & \casevide \\
		Racines et signe d’un polynôme de degré 3 de la forme $a(x-x_1)(x-x_2)(x-x_3)$ & \casevide & \casevide \\
		Résolution d’équations $x^3 = c$ ; racine cubique & \casevide & \casevide \\
		\hline
		\multicolumn{3}{|L{13.5cm}|}{\textbf{Dérivation}} \\
		Nombre dérivé (approche graphique : sécantes, tangente comme position limite) & \casevide & \casevide \\
		Équation réduite de la tangente & \casevide & \casevide \\
		Fonction dérivée de $x \mapsto x^2$, $x \mapsto x^3$ & \casevide & \casevide \\
		Dérivée d’une somme, de $kf$ (k réel), dérivée d’un polynôme de degré $\le 3$ & \casevide & \casevide \\
		Lien entre signe de la dérivée et sens de variation, tableau de variations, extremums & \casevide & \casevide \\
		\hline
		
		% --- Statistiques ---
		\multicolumn{3}{|L{13.5cm}|}{\textbf{Statistiques – Croisement de deux variables catégorielles}} \\
		Tableau croisé d’effectifs, fréquences conditionnelles et marginales & \casevide & \casevide \\
		Compléter un tableau par raisonnement sur les effectifs ou les fréquences conditionnelles & \casevide & \casevide \\
		\hline
		
		% --- Probabilités conditionnelles ---
		\multicolumn{3}{|L{13.5cm}|}{\textbf{Probabilités conditionnelles}} \\
		Probabilité conditionnelle $P_A(B)$, calcul via tableau croisé & \casevide & \casevide \\
		\hline
		\multicolumn{3}{|L{13.5cm}|}{\textbf{Modèle à plusieurs épreuves indépendantes}} \\
		Arbre de probabilités pour deux épreuves indépendantes & \casevide & \casevide \\
		Répétition d’épreuves de Bernoulli identiques et indépendantes (n < 4), calcul de probabilités par arbre & \casevide & \casevide \\
		\hline
		\multicolumn{3}{|L{13.5cm}|}{\textbf{Variable aléatoire – Loi de Bernoulli}} \\
		Variable aléatoire discrète : loi de probabilité, espérance & \casevide & \casevide \\
		Loi de Bernoulli de paramètre $p$, espérance & \casevide & \casevide \\
		Simulation d’échantillons de loi de Bernoulli, observation de la fluctuation d’échantillonnage & \casevide & \casevide \\
		\hline
	\end{longtable}
	
	\section{Calendrier et programmation harmonieuse}
	
	\subsection{Calendrier retenu (zone C, jours ouvrés)}
	Les vacances et jours fériés ont été pris en compte. Les semaines de cours sont les suivantes :
	
	\begin{center}
		\begin{tabular}{|C{3.5cm}|C{5cm}|C{2.5cm}|}
			\hline
			\textbf{Période} & \textbf{Dates} & \textbf{Nb de semaines} \\
			\hline
			Rentrée – Toussaint & 1er sept. – 16 oct. 2026 & 7 \\
			Toussaint – Noël & 2 nov. – 18 déc. 2026 & 7 \\
			Noël – Hiver & 4 janv. – 5 fév. 2027 & 5 \\
			Hiver – Printemps & 22 fév. – 2 avr. 2027 & 6 \\
			Printemps – Fin mai & 19 avr. – 31 mai 2027 & 6 \\
			\hline
			\textbf{Total} & & \textbf{31 semaines} \\
			\hline
		\end{tabular}
	\end{center}
	
	\textit{Jours fériés (tombant un jour de cours) exclus :} 1er janvier 2027 (vendredi), 29 mars 2027 (lundi de Pâques), 6 mai 2027 (Ascension, jeudi), 17 mai 2027 (Pentecôte, lundi). Les 1er et 8 mai 2027 tombent un samedi.
	
	\subsection{Programmation par trimestre (approche spiralée)}
	Les tableaux suivants précisent pour chaque semaine :
	\begin{itemize}
		\item les contenus nouveaux,
		\item les \textbf{points ciblés dans l’évaluation formative} (5 à 10 min en début de séance) incluant systématiquement une question Python,
		\item des remarques utiles.
	\end{itemize}
	Les évaluations formatives mêlent mathématiques et programmation, avec réactivation des notions antérieures.
	
	% --- 1er trimestre ---
	\newpage
	\subsubsection{1\textsuperscript{er} trimestre (septembre – novembre 2026) – 14 semaines}
	
	\rowcolors{2}{gray!8}{white}
	\begin{longtable}{|C{0.7cm}|L{3.8cm}|L{5.5cm}|L{2.3cm}|}
		\hline
		\textbf{Sem.} & \textbf{Contenus nouveaux} & \textbf{Évaluation formative -- points évalués (dont Python)} & \textbf{Remarques} \\
		\hline
		\endhead
		\hline
		\multicolumn{4}{r}{Suite de la page suivante} \\
		\endfoot
		\hline
		\endlastfoot
		
		1-2 & Suites : modes de génération (formule explicite, récurrence), représentation graphique. & Calculs de puissances, racines. Notion de fonction (2nde). \textbf{Python} : boucle \texttt{for} pour calculer les premiers termes d’une suite. & Introduction aux modèles discrets. \\
		\hline
		3-4 & Suites arithmétiques : relation de récurrence, sens de variation, représentation. & \textbf{Spirale :} suites en général (S1-2). \textbf{Python} : fonction qui calcule la somme des termes d’une suite arithmétique (à la main sur petits exemples). & Croissance linéaire. \\
		\hline
		5-6 & Suites géométriques (termes positifs) : relation de récurrence, sens de variation, représentation. & \textbf{Spirale :} suites arithmétiques (S3-4). \textbf{Python} : simulation de l’évolution d’un capital à intérêts composés. & Croissance exponentielle. \\
		\hline
		7-8 & Fonctions polynômes de degré 2 : formes $x^2$, $x^2+b$, $a(x-x_1)(x-x_2)$, racines, signe, extremum. & \textbf{Spirale :} fonctions carré (2nde). \textbf{Python} : tracé de paraboles avec \texttt{matplotlib}. & Optimisation simple. \\
		\hline
		9-10 & Résolution d’équations $x^2=c$. Taux de variation d’une fonction, pente de la sécante. & \textbf{Spirale :} signe d’un polynôme (S7-8). \textbf{Python} : calcul de taux de variation entre deux points. & Préparation à la dérivation. \\
		\hline
		11-12 & Nombre dérivé : approche graphique (sécantes, tangente), équation de la tangente. & \textbf{Spirale :} taux de variation (S9-10). \textbf{Python} : visualisation de la position limite des sécantes avec un curseur. & Introduction à la dérivée. \\
		\hline
		13-14 & Synthèse sur les suites et les fonctions du second degré. Problèmes d’optimisation. & \textbf{Bilan T1 :} suites, polynômes degré 2, tangente. \textbf{Python} : recherche d’extremum par balayage. & Préparation au contrôle. \\
		\hline
		\multicolumn{4}{l}{} \\
		\hline
		\multicolumn{4}{l}{\textbf{Contrôles continus :} semaines 5, 10, 14 (avec exercice Python). \textbf{Formatives :} chaque semaine.} \\
		\hline
	\end{longtable}
	
	% --- 2e trimestre ---
	\newpage
	\subsubsection{2\textsuperscript{e} trimestre (novembre 2026 – février 2027) – 12 semaines}
	
	\rowcolors{2}{gray!8}{white}
	\begin{longtable}{|C{0.7cm}|L{3.8cm}|L{5.5cm}|L{2.3cm}|}
		\hline
		\textbf{Sem.} & \textbf{Contenus nouveaux} & \textbf{Évaluation formative -- points évalués (dont Python)} & \textbf{Remarques} \\
		\hline
		\endhead
		\hline
		\multicolumn{4}{r}{Suite de la page suivante} \\
		\endfoot
		\hline
		\endlastfoot
		
		15-16 & Fonctions polynômes de degré 3 : $ax^3$, $ax^3+b$, racines et signe de $a(x-x_1)(x-x_2)(x-x_3)$. Résolution $x^3=c$. & \textbf{Spirale :} degré 2 (S7-8). \textbf{Python} : tracé de courbes cubiques. & Généralisation. \\
		\hline
		17-18 & Dérivée des fonctions $x^2$, $x^3$, dérivée d’une somme, de $kf$. Dérivée d’un polynôme de degré $\le 3$. & \textbf{Spirale :} nombre dérivé (S11-12). \textbf{Python} : fonction qui calcule la dérivée d’un polynôme (coefficients). & Calcul formel simple. \\
		\hline
		19-20 & Lien entre signe de la dérivée et sens de variation, tableau de variations, extremums. & \textbf{Spirale :} dérivation (S17-18). \textbf{Python} : étude complète d’une fonction (dérivée, tableau). & \textbf{PBL 1 (2h)} : Optimisation d’un volume (ex: boîte sans couvercle). \\
		\hline
		21-22 & Statistiques : tableau croisé d’effectifs, fréquences conditionnelles et marginales. & \textbf{Spirale :} pourcentages (automatismes). \textbf{Python} : traitement de fichiers CSV, filtrage. & Contextes réels. \\
		\hline
		23-24 & Compléter un tableau par raisonnement sur les effectifs ou les fréquences conditionnelles. & \textbf{Spirale :} fréquences conditionnelles (S21-22). \textbf{Python} : calcul de fréquences à partir de listes. & Interprétation. \\
		\hline
		25-26 & Probabilités conditionnelles : $P_A(B)$ via tableau croisé. & \textbf{Spirale :} statistiques (S21-22). \textbf{Python} : simulation de tirages aléatoires avec condition. & Tests diagnostiques. \\
		\hline
		\multicolumn{4}{l}{} \\
		\hline
		\multicolumn{4}{l}{\textbf{Contrôles continus :} semaines 18, 22, 26 (avec exercice Python). \textbf{Formatives :} chaque semaine.} \\
		\hline
	\end{longtable}
	
	% --- 3e trimestre ---
	\newpage
	\subsubsection{3\textsuperscript{e} trimestre (février – mai 2027) – 12 semaines}
	
	\rowcolors{2}{gray!8}{white}
	\begin{longtable}{|C{0.7cm}|L{3.8cm}|L{5.5cm}|L{2.3cm}|}
		\hline
		\textbf{Sem.} & \textbf{Contenus nouveaux} & \textbf{Évaluation formative -- points évalués (dont Python)} & \textbf{Remarques} \\
		\hline
		\endhead
		\hline
		\multicolumn{4}{r}{Suite de la page suivante} \\
		\endfoot
		\hline
		\endlastfoot
		
		27-28 & Modèle à deux épreuves indépendantes : arbre de probabilités, probabilité produit. & \textbf{Spirale :} conditionnelles (S25-26). \textbf{Python} : simulation de deux lancers de dés. & Indépendance. \\
		\hline
		29-30 & Répétition d’épreuves de Bernoulli (n < 4) : arbre, calcul de probabilités. & \textbf{Spirale :} arbres (S27-28). \textbf{Python} : simulation de schéma de Bernoulli pour n=3. & \textbf{PBL 2 (2h)} : Simulation d’un jeu de hasard (ex: pile ou face). \\
		\hline
		31-32 & Variable aléatoire discrète : loi de probabilité, espérance. Loi de Bernoulli, espérance. & \textbf{Spirale :} probabilités (S29-30). \textbf{Python} : calcul d’espérance à partir d’une liste de valeurs. & Interprétation. \\
		\hline
		33-34 & Simulation d’échantillons de loi de Bernoulli, fluctuation d’échantillonnage, écart-type. & \textbf{Spirale :} loi de Bernoulli (S31-32). \textbf{Python} : génération de N échantillons, histogramme des fréquences. & \textbf{PBL 3 (2h)} : Sondage – fluctuation et décision. \\
		\hline
		35-36 & Synthèse probabilités et statistiques. & \textbf{Bilan probabilités/stats}. \textbf{Python} : projet complet de traitement de données. & Préparation au contrôle. \\
		\hline
		37-38 & Révisions et évaluations finales. & \textbf{Bilan général :} tous les thèmes + Python. & Préparation épreuves. \\
		\hline
		\multicolumn{4}{l}{} \\
		\hline
		\multicolumn{4}{l}{\textbf{Contrôles continus :} semaines 30, 34, 38 (avec exercice Python). \textbf{Formatives :} chaque semaine.} \\
		\hline
	\end{longtable}
	
	\section{Évaluations}
	
	\subsection{Évaluations formatives (hebdomadaires)}
	Chaque semaine, une courte activité (5–10 min) mêle mathématiques et Python. Les points évalués sont détaillés dans les tableaux ci-dessus. Ces évaluations ne sont pas notées, mais elles permettent une auto-évaluation et un suivi des compétences algorithmiques.
	
	\subsection{Évaluations sommatives (contrôles continus)}
	Au moins trois contrôles par trimestre, d’une durée de 1 h chacun.  
	Chaque contrôle comporte :
	\begin{itemize}
		\item une partie sans calculatrice (calculs, raisonnement),
		\item une partie avec calculatrice / Python (exercice de programmation : écriture de fonction, complétion de code, simulation, représentation graphique, etc.).
	\end{itemize}
	
	\begin{center}
		\begin{tabular}{|C{3.5cm}|C{3.5cm}|C{3.5cm}|C{3.5cm}|}
			\hline
			\textbf{Trimestre} & \textbf{Contrôle 1} & \textbf{Contrôle 2} & \textbf{Contrôle 3} \\
			\hline
			1\textsuperscript{er} (sem. 5, 10, 14) & Suites (arithm./géom.) \newline + Python (boucles) & Polynômes degré 2, taux de variation \newline + Python (tracés) & Nombre dérivé, tangente \newline + Python (visualisation) \\
			\hline
			2\textsuperscript{e} (sem. 18, 22, 26) & Polynômes degré 3, dérivation (calcul) \newline + Python (dérivée formelle) & Tableau de variations, extremums \newline + Python (étude de fonction) & Statistiques (fréquences conditionnelles) \newline + Python (fichiers) \\
			\hline
			3\textsuperscript{e} (sem. 30, 34, 38) & Probabilités conditionnelles, arbres \newline + Python (simulation) & Loi de Bernoulli, espérance \newline + Python (simulation, histogramme) & Synthèse (probabilités, stats) \newline + Python (projet complet) \\
			\hline
		\end{tabular}
	\end{center}
	
	\subsection{Problèmes basés sur l’apprentissage (PBL) – 2 h chacun}
	Trois PBL sont répartis sur l’année, en séances de 2 h (travail en groupes, compte rendu écrit et code Python rendu) :
	
	\begin{enumerate}
		\item \textbf{PBL 1 (semaine 20)} : \emph{Optimisation d’une boîte sans couvercle}. À partir d’une feuille carrée, découper des carrés aux coins pour former une boîte ; exprimer le volume en fonction de la hauteur, étudier la fonction (dérivée) pour trouver le maximum. \textbf{Python} : tracé de la courbe et recherche de l’extremum par balayage.
		\item \textbf{PBL 2 (semaine 30)} : \emph{Simulation d’un jeu de hasard}. Simuler une expérience aléatoire (par exemple, lancer d’une pièce 3 fois) et comparer les fréquences simulées avec les probabilités théoriques. \textbf{Python} : simulation, arbre de probabilités.
		\item \textbf{PBL 3 (semaine 34)} : \emph{Sondage – fluctuation d’échantillonnage}. Simuler N échantillons de taille n d’une loi de Bernoulli, représenter l’histogramme des fréquences, observer l’écart-type, discuter de la précision d’un sondage. \textbf{Python} : génération d’échantillons, calcul d’écart-type.
	\end{enumerate}
	
	\section{Devoirs maison (DM) -- progression, Python et utilisation de l’IA}
	
	Afin de favoriser une acquisition à long terme, un devoir maison est distribué \textbf{toutes les deux semaines} (18 DM sur l’année).  
	Chaque DM porte sur l’ensemble du programme traité depuis le 1er septembre, et comporte :
	\begin{itemize}
		\item des exercices de calcul et de raisonnement,
		\item un exercice de programmation Python,
		\item une mission spécifique d’utilisation de l’IA (consultation, débuggage, confrontation de méthodes ou critique de réponse).
	\end{itemize}
	
	L’élève doit joindre à sa copie un \textbf{journal de bord} retraçant ses échanges avec l’IA et les décisions prises à la suite de ces échanges.
	
	\begin{center}
		\small
		\rowcolors{2}{gray!8}{white}
		\begin{longtable}{|C{0.7cm}|C{0.9cm}|L{6.0cm}|L{6.0cm}|}
			\hline
			\textbf{DM} & \textbf{Sem.} & \textbf{Thèmes évalués (mathématiques + Python)} & \textbf{Mission IA spécifique} \\
			\hline
			\endhead
			\hline
			\multicolumn{4}{r}{Suite à la page suivante} \\
			\endfoot
			\hline
			\endlastfoot
			
			1 & 3 & Suites : modes de génération, représentation. Python : calcul de termes. & Demander à l’IA d’expliquer la différence entre une suite arithmétique et une suite géométrique avec un exemple concret. \\
			\hline
			2 & 5 & Suites arithmétiques : relation de récurrence, sens de variation. Python : somme. & Utiliser l’IA pour générer un problème de placement simple et le résoudre. \\
			\hline
			3 & 7 & Suites géométriques : relation, sens de variation. Python : simulation d’évolution. & Demander à l’IA d’illustrer la croissance exponentielle avec un exemple de population. \\
			\hline
			4 & 9 & Polynômes degré 2 : forme factorisée, racines, signe. Python : tracé. & Utiliser l’IA pour vérifier le signe d’une expression et justifier. \\
			\hline
			5 & 11 & Taux de variation, pente de la sécante. Python : calcul de taux. & Demander à l’IA d’expliquer la notion de taux de variation avec une interprétation concrète (vitesse). \\
			\hline
			6 & 13 & Nombre dérivé, tangente. Python : visualisation des sécantes. & Utiliser l’IA pour démontrer que la tangente est la position limite des sécantes. \\
			\hline
			7 & 16 & Polynômes degré 3 : racines, signe. Python : tracé. & Demander à l’IA de proposer une factorisation d’un polynôme du 3e degré dont une racine est connue. \\
			\hline
			8 & 18 & Dérivée des polynômes degré $\le 3$. Python : fonction dérivée. & Utiliser l’IA pour vérifier les règles de dérivation et les appliquer. \\
			\hline
			9 & 20 & Lien dérivée-variations, tableau de variations. Python : étude complète. & Demander à l’IA de critiquer un tableau de variations et de proposer des améliorations. \\
			\hline
			10 & 22 & Statistiques : tableau croisé, fréquences conditionnelles. Python : traitement de fichier. & Utiliser l’IA pour interpréter des fréquences conditionnelles issues d’un tableau. \\
			\hline
			11 & 24 & Compléter un tableau croisé. Python : filtrage de données. & Demander à l’IA de proposer des questions pertinentes à partir d’un tableau. \\
			\hline
			12 & 26 & Probabilités conditionnelles via tableau. Python : simulation de tirage. & Utiliser l’IA pour expliquer la notion de probabilité conditionnelle avec un exemple médical. \\
			\hline
			13 & 28 & Arbre de probabilités pour deux épreuves indépendantes. Python : simulation. & Demander à l’IA de construire un arbre à partir d’un énoncé et de vérifier les calculs. \\
			\hline
			14 & 30 & Répétition de Bernoulli (n<4). Python : simulation de schéma de Bernoulli. & Utiliser l’IA pour dénombrer les chemins dans un arbre et expliquer la formule. \\
			\hline
			15 & 32 & Loi de Bernoulli, espérance. Python : calcul d’espérance. & Demander à l’IA d’interpréter l’espérance dans un jeu de hasard. \\
			\hline
			16 & 34 & Simulation d’échantillons, fluctuation. Python : histogramme des fréquences. & Utiliser l’IA pour discuter de la précision d’un sondage à partir des simulations. \\
			\hline
			17 & 36 & Synthèse probabilités et statistiques. Python : projet complet. & Demander à l’IA un retour sur la structure du code et proposer des améliorations. \\
			\hline
			18 & 38 & Bilan général. Python : projet complet. & Demander à l’IA de suggérer des extensions du problème traité. \\
			\hline
		\end{longtable}
	\end{center}
	
	\textbf{Remarques :}
	\begin{itemize}
		\item Les DM sont à rendre une semaine après la distribution.
		\item La mission IA doit être clairement identifiée dans le journal de bord (copie d’écran ou retranscription des questions/réponses, suivie de l’analyse personnelle de l’élève).
		\item La correction est effectuée en classe, avec une attention particulière aux erreurs fréquentes, aux méthodes de résolution et aux bonnes pratiques d’utilisation de l’IA.
	\end{itemize}
	
	\section{Conseils pour la mise en œuvre}
	\begin{itemize}
		\item \textbf{Rituels} : consacrer 5 à 10 minutes par séance à des exercices de calcul (mental, littéral) et à de courts extraits de code Python (lecture, complétion, débogage).
		\item \textbf{Travail en îlots} : favoriser les échanges entre élèves lors des phases de recherche et des PBL.
		\item \textbf{Numérique} : utiliser Python régulièrement (au moins une séance sur deux). Insister sur la programmation modulaire (fonctions) et la documentation du code.
		\item \textbf{IA éthique} : former les élèves à formuler des requêtes précises et à évaluer de manière critique les réponses de l’IA (vérification des calculs, cohérence, limites). Insister sur le fait que l’IA est un \emph{outil d’aide}, pas un substitut au raisonnement personnel.
		\item \textbf{Différenciation} : proposer des exercices de difficulté progressive et des approfondissements pour les élèves plus à l’aise (notamment en programmation et en interaction avec l’IA).
		\item \textbf{Trace écrite} : veiller à ce que chaque élève dispose d’un cahier de cours structuré (définitions, propriétés, exemples) et d’un cahier de programmation regroupant les fonctions et scripts clés.
	\end{itemize}
	
	\section*{Conclusion}
	Cette programmation respecte les 31 semaines de cours disponibles et couvre l’intégralité du programme de première technologique (tronc commun), avec une intégration systématique de Python et de l’utilisation critique de l’IA dans les apprentissages et les évaluations. La progressivité spiralaire, appliquée aussi aux compétences algorithmiques et numériques, garantit que les notions mathématiques, la programmation et le jugement critique sont solidement maîtrisés en fin d’année. Les évaluations variées (formatives, sommatives, PBL, DM) permettent de suivre la progression de chaque élève et de développer ses compétences de modélisation, de raisonnement, de communication, de codage et de collaboration avec les outils d’intelligence artificielle.
	
\end{document}